A Hawkes process is a counting process whose intensity is driven by its own past.
Each event raises the intensity of future events by an amount that depends on the two types
involved and decays with the time elapsed.

\begin{defi}\label{def.hawkesprocess}
 A $\numEventTypes$-dimensional counting process $\multiCountingProc$ is called a Hawkes process
 if it admits an absolutely continuous compensator $\compensator$ with intensities
 \begin{equation}\label{eq.intensity_ordHawkes}
  \intensity(t) = \baseIntensity_e
  + \sum_{\eone=1}^{\numEventTypes}\intzerot \hawkesKernel\subscriptee(t-s) \, d\countingProc[\eone](s),
  \qquad e = 1,\dots,\numEventTypes,
 \end{equation}
 for some non-negative baselines $\baseIntensity_e \geq 0$,
 and some non-negative locally integrable functions $\hawkesKernel\subscriptee \geq 0$
 supported on the non-negative half line.
\end{defi}

The matrix-valued function
$t\mapsto [\hawkesKernel\subscriptee(t)]_{e,\eone = 1, \dots, \numEventTypes}$
is referred to as the kernel of the Hawkes process $\multiCountingProc$.
If all the kernel functions are integrable, we write
\begin{equation}\label{eq.branchingMatrix}
 \branchingMatrix\subscriptee := \Vert \hawkesKernel\subscriptee \Vert_{\Lone}
 = \int_{0}^{\infty} \hawkesKernel\subscriptee(t) \, dt
\end{equation}
for the matrix of $\Lone$ norms,
and we call the spectral radius $\branchingRatio$ of $\branchingMatrix$ the branching ratio of
the process;
if some of the kernel functions are not integrable, the branching ratio is set to $+\infty$.
Every statement of this chapter about stability is a statement about $\branchingRatio$,
and Section \ref{sec.stability} is devoted to reading it.
For the moment we record only that a $\numEventTypes$-dimensional Hawkes process admits a
stationary version with finite mean intensity if and only if $\branchingRatio < 1$.

\begin{remark}[The order of the indices]\label{remark.indexOrder}
 In \eqref{eq.intensity_ordHawkes},
 $\hawkesKernel\subscriptee$ is the influence that an event of type $\eone$
 exerts on an event of type $e$:
 the row is excited and the column excites.
 Consequently the \emph{column} sum $\sum_{e}\branchingMatrix\subscriptee$
 is the expected number of direct offspring of a single event of type $\eone$,
 and the products $\intensity[ ] = \baseIntensity + \excitation\decayedCounts$
 and $\stationaryIntensity = (I - \branchingMatrix)^{-1}\baseIntensity$
 are products on column vectors.
 Since a matrix and its transpose share a spectral radius,
 a kernel transposed by mistake passes every stability test of
 Section \ref{sec.stability} and produces a path of the right total rate;
 the convention is the only defence against it.
\end{remark}

\subsubsection{The cluster representation}
\label{sec.clusterRepresentation}
\begin{defi}\label{def.galtonWatson}
 Let $\Gamma \geq 0$ be a $\numEventTypes \times \numEventTypes$ matrix.
 A \emph{multitype Galton--Watson process} with mean offspring matrix $\Gamma$
 is a population in which every individual of type $\eone$,
 independently of every other,
 produces a random number of children of type $e$ of mean
 $\Gamma\subscriptee$, for each $e$,
 and the generations are counted one after the other.
\end{defi}
With the index convention of Remark \ref{remark.indexOrder},
the mean offspring matrix of Definition \ref{def.galtonWatson} is
$\branchingMatrix$ itself;
part of the literature writes the transpose,
and then the \emph{row} sums are the offspring counts.

Definition \ref{def.hawkesprocess} describes the process through its intensity.
There is a second description, due to \cite{HO74clu},
which describes the same process through a branching construction,
and which is the one that explains why the order flow carries information at all.

\begin{prop}[\cite{HO74clu}]\label{prop.clusterRepresentation}
 Assume $\hawkesKernel \geq 0$ and $\branchingRatio < 1$.
 Construct a point process as follows.
 \emph{i.} Immigrants of type $e$ arrive as a homogeneous Poisson process of rate
 $\baseIntensity_e$, independently across $e$.
 \emph{ii.} Every event of type $\eone$, immigrant or not,
 independently generates its direct offspring of type $e$
 as an inhomogeneous Poisson process on $(t,\infty)$ of intensity $\hawkesKernel\subscriptee(s-t)$,
 so that the number of such offspring is Poisson with mean $\branchingMatrix\subscriptee$.
 Then the superposition of all events so constructed is a stationary Hawkes process
 with baselines $\baseIntensity$ and kernel $\hawkesKernel$,
 and conversely every stationary Hawkes process with finite intensity admits this representation.
\end{prop}

The descendants of a single immigrant form a multitype Galton--Watson process
whose mean offspring matrix is $\branchingMatrix$,
and the population dies out almost surely, with finite expected total size,
precisely when $\branchingRatio < 1$.
The events of the process therefore partition into \emph{clusters},
each consisting of one immigrant and all of its descendants.


\begin{defi}[Exponential-kernel Hawkes process]\label{def.exponentialHawkes}
 A $\numEventTypes$-dimensional Hawkes process $\multiCountingProc$ is called an
 \emph{exponential-kernel Hawkes process} if its kernels are
 \begin{equation}\label{eq.exponentialKernel}
  \hawkesKernel\subscriptee(t) = \excitation\subscriptee \, e^{-\decay t},
  \qquad t \geq 0,
 \end{equation}
 with $\excitation\subscriptee \geq 0$ and a single decay rate $\decay > 0$
 shared by every pair of types.
 We call $\excitation$ the excitation matrix and $\decay$ the decay.
 For each type $e$ we call
 \begin{equation}\label{eq.decayedCounts}
  \decayedCounts_e(t) := \sum_{j \, : \, \arrivalTimes_{j} < t}
  e^{-\decay (t - \arrivalTimes_{j})}
 \end{equation}
 the decayed counts.
\end{defi}

From \eqref{eq.branchingMatrix},
\begin{equation}\label{eq.GammaIsAOverBeta}
 \branchingMatrix = \frac{\excitation}{\decay},
\end{equation}
so the branching matrix is dimensionless while the excitation matrix is a rate.

\begin{remark}\label{remark.countsAgainstRates}
 Notice that $\sum_e \branchingMatrix\subscriptee$ and $\sum_e \excitation\subscriptee$
 are different quantities, differing by $\decay$:
 the first is the expected number of direct offspring of one event of type $\eone$,
 the second is the jump in the total intensity that the same event produces,
 measured per unit of time.
\end{remark}

The inequality under the sum in \eqref{eq.decayedCounts} is strict.
It is what makes $\decayedCounts$ left-continuous,
and therefore makes
\begin{equation}\label{eq.intensityFromState}
 \intensity[ ](t) = \baseIntensity + \excitation \decayedCounts(t)
\end{equation}
predictable, as Definition \ref{def.compensator} requires of an intensity.

Between events $\decayedCounts$ solves $d\decayedCounts = -\decay \decayedCounts \, dt$,
and at an event of type $e$ it gains the unit vector $\unitVector_e$:
\begin{equation}\label{eq.stateDynamics}
 d\decayedCounts(t) = -\decay \decayedCounts(t) \, dt + d\multiCountingProc(t).
\end{equation}

\begin{thm}\label{thm.markovState}
 Let $\excitation \geq 0$, $\decay > 0$ and $\baseIntensity \geq 0$.
 Then $\decayedCounts$ is a piecewise-deterministic Markov process on
 $\R_{+}^{\numEventTypes}$,
 with flow $\varphi_t(z) = e^{-\decay t}z$,
 jump rate $\totalIntensity(z) = \mathbf{1}^{\transpose}(\baseIntensity + \excitation z)$
 and jump kernel
 $\sum_{e} \frac{\intensity(z)}{\totalIntensity(z)} \delta_{z + \unitVector_e}$,
 and its extended generator is
 \begin{equation}\label{eq.generator}
  \generatorL f(z)
  = -\decay \, z \cdot \gradient f(z)
  + \sum_{e=1}^{\numEventTypes} \intensity(z)\big( f(z + \unitVector_e) - f(z) \big),
  \qquad \intensity(z) = \baseIntensity_e + (\excitation z)_e ,
 \end{equation}
 on every $f \in C^{1}(\R^{\numEventTypes})$ such that
 $\Expectation \int_{0}^{t} \sum_{e} \intensity(\decayedCounts(s))
 \lvert f(\decayedCounts(s) + \unitVector_e) - f(\decayedCounts(s)) \rvert \, ds < \infty$
 for every $t > 0$.
 The pair $(\multiCountingProc, \decayedCounts)$ is then Markov as well,
 with generator
 $\generatorL f(n,z) = -\decay \, z \cdot \gradient_{z} f(n,z)
 + \sum_{e} \intensity(z)\big( f(n + \unitVector_e, z + \unitVector_e) - f(n,z) \big)$.
\end{thm}
\begin{proof}
 The intensity \eqref{eq.intensityFromState} and the jump kernel are functions of
 $\decayedCounts$ alone, so $\decayedCounts$ is a Markov process and
 $\multiCountingProc$ is the process counting its jumps by type.
 Between jumps $\decayedCounts$ follows the flow, and
 $\frac{d}{dt} f(\varphi_t(z)) = -\decay \, \varphi_t(z)\cdot \gradient f(\varphi_t(z))$.
 By Definition \ref{def.compensator} the jump measure of $\decayedCounts$ has compensator
 $\sum_{e} \intensity(\decayedCounts(s)) \, \delta_{\decayedCounts(s) + \unitVector_e} \, ds$,
 so the sum of the jumps of $f(\decayedCounts)$ minus the integral of the second term of
 \eqref{eq.generator} is a local martingale.
 Adding the two and localising gives \eqref{eq.generator} up to localisation,
 which the stated integrability removes.
\end{proof}
For the general statement on piecewise-deterministic Markov processes and for the domain of
the extended generator, see \cite[Theorem 5.5]{Dav84pie};
the Markov property of the exponential-kernel case is due to \cite{Oak75mar}.

\begin{remark}\label{remark.noSubcriticality}
 Notice that Theorem \ref{thm.markovState} assumes nothing about $\branchingRatio$.
 The integrability it requires holds on every finite horizon at every $\branchingRatio$,
 both for $f(z) = z$ and for $f(z) = z z^{\transpose}$,
 the jump increments of the latter being of size $O(\lvert z \rvert)$.
 What subcriticality buys is stated in Section \ref{sec.stability},
 and it is not this.
\end{remark}

\begin{remark}[choosing the kernel is choosing the size of the state]\label{remark.stateSize}
 A single $\decay$ shared by all pairs buys a state of dimension $\numEventTypes$.
 A decay $\decay\subscriptee$ per pair would require one decayed count per pair,
 a state of dimension $\numEventTypes^2$,
 and every expression of Sections \ref{sec.stability} and \ref{sec.secondOrder} would have to
 be rewritten on it.
 A decay per \emph{type} keeps the state of dimension $\numEventTypes$
 but loses \eqref{eq.totalIntensityBetweenEvents}, and with it the scheme of
 Theorem \ref{thm.dassiosZhao}.
 What the single decay costs is that the model has one excitation timescale
 and cannot describe a market whose market orders relax quickly and whose cancellations
 relax slowly.
\end{remark}

From now on, unless otherwise stated,
every Hawkes process is an exponential-kernel Hawkes process.

Integrating \eqref{eq.stateDynamics} over $[0,\timeHorizon]$ with
$\decayedCounts(0) = \multiCountingProc(0) = 0$ gives
$\int_{0}^{\timeHorizon} \decayedCounts_{\eone} =
\big( \countingProc[\eone](\timeHorizon-) - \decayedCounts_{\eone}(\timeHorizon) \big) / \decay$,
and therefore
\begin{equation}\label{eq.compensatorClosedForm}
 \compensator_e(\timeHorizon)
 = \baseIntensity_e \timeHorizon
 + \sum_{\eone=1}^{\numEventTypes} \frac{\excitation\subscriptee}{\decay}
 \Big( \countingProc[\eone](\timeHorizon-) - \decayedCounts_{\eone}(\timeHorizon) \Big).
\end{equation}
The compensator is thus available exactly, at any time, from the state and the counts,
with no numerical integration.

\begin{remark}\label{remark.readPreJump}
 Both terms of \eqref{eq.compensatorClosedForm} are read immediately before $\timeHorizon$:
 $\countingProc[\eone]$ is right-continuous and $\decayedCounts_{\eone}$ is left-continuous,
 so $\countingProc[\eone] - \decayedCounts_{\eone}$ jumps by one at every arrival of type
 $\eone$ and the two readings do not agree there.
 They differ only at arrival times, which is precisely where
 Section \ref{sec.simulation} evaluates the compensator:
 with $\decay = 1$, a single arrival at $t = 1/2$ and $\timeHorizon = 1/2$,
 the integral is $0$ while the post-jump reading gives $1$.
 It is also why we compute the compensator independently of the simulator rather than
 accumulating it along the path:
 an independent computation is a test, an accumulated one is a tautology.
\end{remark}
